% §11.5 Bridge
%
% Monograph register: recap (with the corrected axis assignment, RP = FPR-zero,
% coRP = FNR-zero) and forward to chapters 12-13. Cross-refs
% sec:miller-rabin-formal, thm:miller-rabin-formal, sec:las-vegas-vs-monte-carlo,
% def:monte-carlo, sec:rp-bpp, def:rp, def:perfect-filter-line,
% sec:one-sided-two-sided, sec:sample-complexity, thm:chernoff,
% sec:binomial-estimator, ch:approximate-relations, ch:algebra-types. DEFINED:
% sec:ch11-bridge.

\section{Bridge}
\label{sec:ch11-bridge}

The chapter traced one figure from four angles.
\Cref{sec:miller-rabin-formal} closed the chapter-3 deferral, making
Miller-Rabin a formal $\Bernoulli{\BoolType}$ with $\fnrate = 0$ and
$\fprate \le 1/4$ per round (\Cref{thm:miller-rabin-formal}).
\Cref{sec:las-vegas-vs-monte-carlo} separated the Monte Carlo regime, where
$\Bernoulli{\BoolType}$ analysis applies, from its exact-output Las Vegas
sibling. \Cref{sec:rp-bpp} read the Monte Carlo classes as locations on the
rate plane of \Cref{def:perfect-filter-line}: RP (\Cref{def:rp}) on the
\emph{FPR-zero} axis (certain accept), coRP on the \emph{FNR-zero} axis (certain
reject, the Bloom-filter shape), ZPP at the origin, BPP in the interior, with
the one-sided-versus-two-sided split of \Cref{sec:one-sided-two-sided}
separating RP and coRP from BPP. \Cref{sec:sample-complexity} quantified the
repetition count, the Chernoff bound (\Cref{thm:chernoff}) giving the
$k = \Theta(\log(1/\delta))$ that makes randomized algorithms practical and
sharpening the chapter-7 variance interval (\Cref{sec:binomial-estimator}) from
polynomial to exponential decay.

The picture extends in two directions before Part~III closes.
\Cref{ch:approximate-relations} lifts the analysis to $\Bernoulli{R}$ for binary
relations, where the question is whether $x$ and $y$ stand in a relation rather
than whether $x$ belongs to a set, the setting for joins and pair predicates; a
Monte Carlo relation test is a $\Bernoulli{R}$ instance with a per-pair rate, and
amplification carries over unchanged. \Cref{ch:algebra-types} closes Part~III by
giving the type algebra under which the chapter-10 sketches, this chapter's
randomized algorithms, and the chapter-12 relations are all instances of a few
constructors over $\Bernoulli{T}$, composing by structure.
